% 数学建模论文LaTeX模板 - 中文版
% 使用XeLaTeX编译

\documentclass[12pt,a4paper]{article}

% ============ 中文支持 ============
\usepackage[UTF8]{ctex}

% ============ 数学公式 ============
\usepackage{amsmath,amssymb,amsfonts}

% ============ 图表 ============
\usepackage{graphicx}
\usepackage{float}
\usepackage{booktabs}
\usepackage{multirow}
\usepackage{subcaption}

% ============ 页面设置 ============
\usepackage{geometry}
\geometry{
    top=2.5cm,
    bottom=2.5cm,
    left=2.5cm,
    right=2.5cm
}

% ============ 页眉页脚 ============
\usepackage{fancyhdr}
\pagestyle{fancy}
\fancyhf{}
\fancyhead[C]{\small 论文标题}
\fancyfoot[C]{\thepage}
\renewcommand{\headrulewidth}{0.4pt}
\setlength{\headheight}{14pt}

% ============ 超链接 ============
\usepackage{hyperref}
\hypersetup{
    colorlinks=true,
    linkcolor=blue,
    citecolor=red,
    urlcolor=blue
}

% ============ 代码高亮 ============
\usepackage{listings}
\usepackage{xcolor}
\lstset{
    basicstyle=\small\ttfamily,
    keywordstyle=\color{blue},
    commentstyle=\color{green!60!black},
    stringstyle=\color{red},
    numbers=left,
    numberstyle=\tiny\color{gray},
    frame=single,
    breaklines=true,
    backgroundcolor=\color{gray!10},
    tabsize=4
}

% ============ 标题信息 ============
\title{\textbf{论文标题}}
\author{
    队员1\thanks{建模手} \quad
    队员2\thanks{编程手} \quad
    队员3\thanks{写作手} \\[0.5em]
    \small 学校名称
}
\date{\today}

\begin{document}

% ============ 封面 ============
\maketitle
\thispagestyle{empty}

% ============ 摘要 ============
\begin{abstract}
本文研究了[问题]。针对问题一，建立了[模型名]，采用[方法]求解，得到[结果]。针对问题二，构建了[模型名]，利用[方法]求得最优解为[数值]。灵敏度分析表明模型具有良好的鲁棒性。

\bigskip
\noindent\textbf{关键词：} 关键词1；关键词2；关键词3；关键词4
\end{abstract}

\newpage
\tableofcontents
\newpage

% ============ 一、问题重述 ============
\section{问题重述}

\subsection{问题背景}
[用自己的语言描述问题背景]

\subsection{问题要求}
[列出问题要求]

% ============ 二、问题分析 ============
\section{问题分析}

\subsection{问题一分析}
[分析问题一]

\subsection{问题二分析}
[分析问题二]

% ============ 三、模型假设 ============
\section{模型假设}

\begin{enumerate}
    \item 假设1：[内容]
    \item 假设2：[内容]
\end{enumerate}

% ============ 四、符号说明 ============
\section{符号说明}

\begin{table}[H]
\centering
\caption{符号说明}
\begin{tabular}{cll}
\toprule
\textbf{符号} & \textbf{含义} & \textbf{单位} \\
\midrule
$x$ & 变量 & 单位 \\
\bottomrule
\end{tabular}
\end{table}

% ============ 五、模型建立与求解 ============
\section{模型建立与求解}

\subsection{问题一模型}

\subsubsection{模型建立}
[模型建立过程]

\subsubsection{模型求解}
[求解过程]

\subsubsection{结果分析}
[结果分析]

\subsection{问题二模型}
[类似结构]

% ============ 六、灵敏度分析 ============
\section{灵敏度分析}

[灵敏度分析]

% ============ 七、模型评价与推广 ============
\section{模型评价与推广}

\subsection{模型优点}
\begin{enumerate}
    \item 优点1
    \item 优点2
\end{enumerate}

\subsection{模型缺点}
\begin{enumerate}
    \item 缺点1
    \item 缺点2
\end{enumerate}

\subsection{改进方向}
\begin{enumerate}
    \item 改进1
    \item 改进2
\end{enumerate}

% ============ 参考文献 ============
\begin{thebibliography}{99}

\bibitem{ref1} 作者. 标题[J]. 期刊, 年份.

\end{thebibliography}

% ============ 附录 ============
\appendix
\section{附录A：代码}

\begin{lstlisting}[language=Python, caption={主要代码}]
# Python代码
\end{lstlisting}

\end{document}
